\documentclass[10pt,a4paper]{amsart}
\usepackage[margin=19mm]{geometry}
\usepackage{amsmath,amssymb,mathtools}
\usepackage[scr=rsfs,cal=euler]{mathalfa}
\usepackage{xfrac}
\usepackage[unicode,hidelinks,bookmarks=false]{hyperref}
\newcommand{\ie}{i.e.\ }
\newcommand{\cf}{cf.\ }
\newcommand{\resp}{respectively\ }
\newcommand{\Subsection}{\subsection}
\providecommand{\nobreakdash}{\nobreak\hbox{-}}
\setlength{\emergencystretch}{2em}
\multlinegap0pt
\usepackage[shortlabels]{enumitem}
\usepackage{bm}
\usepackage{xcolor}
\usepackage{float}
\newtheorem{theorem}{Theorem}
\newtheorem{lemma}{Lemma}
\usepackage{cleveref}
\crefname{theorem}{Theorem}{Theorems}
\crefname{lemma}{Lemma}{Lemmas}
\newcommand{\R}{\mathbb{R}}
\renewcommand{\S}{\mathcal{S}}
\DeclareMathOperator{\Vol}{Vol}
\newcommand{\density}{\delta}
\title{Multiplication tables for integers with restricted prime factors}
\author{Jeremy Schlitt}
\date{Source: arXiv:2603.19212v2 (CC BY 4.0)}
\begin{document}
\maketitle

\noindent\emph{Source and licence.} Multiplication tables for integers with restricted prime factors. Version arXiv:2603.19212v2 (CC BY 4.0), distributed by arXiv under the Creative Commons Attribution 4.0 International licence, http://creativecommons.org/licenses/by/4.0/, as recorded in the arXiv licence metadata for this exact version and shown on the abstract page https://arxiv.org/abs/2603.19212v2. Full licence notice: \texttt{licenses/arxiv-2603.19212-CC-BY-4.0.txt}.\par\smallskip

\noindent\textit{Editorial scope.} Counted unit: the article's lemma Yk (source0.tex lines 503--511). The article's complete original proof of it is delivered with it (source0.tex lines 1090--1162, one \texttt{proof} environment of 73 lines, which the article places away from the statement). No mathematics is written, repaired or re-derived: the statement and the proof are the article's own lines, in the article's own notation, and no retained line is substituted or re-flowed.  Retained uncounted as the source-local context the proof needs: lines 142--144; lines 152--168; lines 460--463; lines 486--490; lines 498--500.  Nothing was omitted from any retained span, so the package carries no omission record.  No retained line contains a citation, so the wrapper carries no bibliography and no bibliography entry.  No retained line contains a figure environment, an image-inclusion command or drawing code, so no image is delivered and the package declares no asset dependency.  Editorial formatting only, in addition: an AMS \texttt{amsart} preamble that restates the article's own declarations and macros used by the retained lines, namely the environments \texttt{theorem}, \texttt{lemma} on separate counters, together with the standard packages those lines need.  The retained environments print in this document as Theorem 1, Lemma 1 in order of appearance, whereas the article prints the same items with the numbering of its own full document.  The complete native source of the article as distributed in the arXiv e-print for this exact version is bundled unchanged as \texttt{sources/196761/source0.tex}.
	\begin{equation}\label{definition of delta}
		\Big|\#\big(Q\cap[1,x]\big) - \frac{ \delta x}{\log x}\Big|\le \frac{\kappa x}{(\log x)^2}, \quad ( x\ge2 ).
	\end{equation}

	\begin{theorem} \label{mainthm}
		Let $c \in (0,1/10]$ and $\kappa \ge 2$ be given. Let $y_0=y_0(c,\kappa)$ be a sufficiently large number in terms of $c,\kappa$. Suppose that ${y_0 \leq y \leq \sqrt{x}}$. One has, uniformly over all sets of primes $Q$ satisfying \eqref{definition of delta} with $\delta \ge c$, that
		\begin{equation}
			H_Q(x,y,2y) \asymp_{\kappa,c} \frac{x}{(\log x)^{1-\density}} \cdot \frac{1}{(\log y)^{G(\density)}} \cdot E(y;\density),
		\end{equation}
		where $$G(\density) := \begin{cases}
			1-\density, \quad &\density \leq \frac{1}{\log 4};\\
			\noalign{\vskip-12pt}\\
			\density -  \frac{1+\log(\density \log 2)}{\log 2}, & \density \geq \frac{1}{\log 4},
		\end{cases}$$
		and
		$$E(y;\density) := \begin{cases}
			\max\big\{\frac{1}{\log 4} - \density, \frac{1}{\sqrt{\log \log y}}\big\} & \text{if } \density \leq \frac{1}{\log 4}, \\
			\noalign{\vskip 5pt}
			\frac{1}{(\log \log y)^{3/2} \max\big\{(\density - \frac{1}{\log 4})^2, \frac{1}{\log \log y}\big\}} & \text{if } \density > \frac{1}{\log 4}.
		\end{cases}$$
	\end{theorem}

	\begin{lemma}\label{lem: size of lambda}
		Let $c,\kappa$ be as in \Cref{mainthm}. Then there exists a constant $R =R(c,\kappa)$ depending only on $c,\kappa$ such that uniformly over $j \geq 1$, one has
		$$2^{j-R} \leq \log(\lambda_j) \leq 2^{j+R}.$$
	\end{lemma}

	\begin{lemma}\label{LB tau} Let $c,\kappa$ be as in \Cref{mainthm}. Let $M = M(c,\kappa)$ be a sufficiently large constant depending only on $c,\kappa$. Suppose that $b_j = 0$ for all $j < M$, and $b_j \leq M j$ for all $j$. One has
		\begin{equation}
			\sum_{a \in \mathscr{A}(\bm{b})} \frac{\tau(a)}{a} \gg_{c,\kappa} \frac{(2 \density \log 2)^k}{b_M! \cdots b_h!}.
		\end{equation}
	\end{lemma}

	\begin{equation}\label{L to factorials}
		\sum_{\substack{a \in \mathscr{A}(\bm{b})}} \frac{L(a)}{a} \gg_{c,\kappa}  \frac{(2 \density \log 2)^k}{b_M! \cdots b_h!} \bigg[ 1 + \sum_{j =M}^h 2^{-j +b_M + \cdots +b_j} \bigg]^{-1}.
	\end{equation}

	\begin{lemma}\label{lem from L to Yk}Let $c,\kappa$ be as in \Cref{mainthm}, and let $M=M(c,\kappa)$ be as in \Cref{LB tau}. Define $$\Tilde{v}:=v-2M=\Big\lfloor\frac{\log_2 y}{\log 2}\Big\rfloor - 2M,$$ and suppose $M+1 \leq k \leq v.$ Fix $C > 2$. Then for some set $\mathscr{A}_k \subseteq \S_Q$ of squarefree integers $a \leq y^{1/8}$ satisfying $\omega(a) = k$, we have
		$$\sum_{\substack{a \in \mathscr{A}_k}} \frac{L(a)}{a} \gg_{c,\kappa} (2v \density \log 2)^k \frac{\Vol(\mathcal{Y}_k(\Tilde{v},C))}{2^C}, $$
		where $\mathcal{Y}_k(\Tilde{v},C)$ is defined to be the set of $\bm{\xi} \in \mathbb{R}^k$ such that
		\begin{enumerate}[label=(\roman*)]
			\item $0 \leq \xi_1 \leq \cdots \leq \xi_k \leq 1$;
			\item $\xi_{M+i^2} > \frac{i}{\Tilde{v}}$ and $\xi_{k+1-(M+i^2)} < 1-\frac{i}{\Tilde{v}}$ for $1 \leq i \leq \sqrt{k-M}$;
			\item $\Tilde{v} \xi_i \geq \max\{i-1, i+ (\min\{i,k-i\})^{1/7}-C\}.$
		\end{enumerate}    
	\end{lemma}

\begin{proof}[Proof of \Cref{lem from L to Yk}]
	Let $h = \Tilde{v}+M-1$. Suppose that $\bm{b} = (b_1,\dots,b_h)$ is such that $b_1 +\cdots + b_h = k$ and $b_j = 0$ for $j < M$. Further assume that
	\begin{equation}\label{eq 1 lem 3.5}
		b_j \leq 2^{j/10}\text{ and } b_{h+1-j} \leq 2^{(M+j)/10} \text{ for } 1 \leq j \leq h.
	\end{equation}
	We claim that any integer $a \in \mathscr{A}(\bm{b})$ with $\bm{b}$ as above satisfies $a \leq y^{1/8}$ as long as $M$ is sufficiently large. To see this, note that by \Cref{lem: size of lambda}, there exists $R = R(c,\kappa) > 0$ such that 
	$$\log a = \sum_{p|a} \log p = \sum_{j=1}^h \sum_{\substack{p|a\\ p \in D_j}} \log p \le \sum_{j=1}^h b_j 2^{j+R} ,$$
	as there are $b_j$ primes dividing $a$ in each interval $D_j$, and each prime $p \in D_j$ satisfies $ \log p \leq  \exp[2^{j+R}]$. By the assumption \eqref{eq 1 lem 3.5}, we have 
	\begin{align*}
		\sum_{j=1}^h b_j 2^{j+R} &\leq \sum_{j = 1}^h 2^{(9j+M+h+1)/10+R}\\
		&\leq 2^{(h+M+1)/10 + R} \sum_{j = 1}^h2^{9(j-h)/10}\\
		&\leq 2^{h+M/10+R+2}\\ & \leq \frac{\log y}{2^{9M/10 - R -2}}.
	\end{align*}
	By choosing $M$ sufficiently large in terms of $R$ (hence large in terms of $c,\kappa$), we end up with $\log a \leq (\log y)/8$, hence $a \leq y^{1/8}$. 
	
	At this point, we relabel our $b_i$'s, writing $g_i = b_{M+i-1},\;(1 \leq i \leq \Tilde{v})$. Suppose that 
	\begin{equation}\label{eq 2 lem 3.5}
		g_i \leq M+i^2\text{ and } g_{\Tilde{v}-i+1} \leq M+i^2 \text{ for } 1 \leq i \leq \Tilde{v}.
	\end{equation}
	This condition implies condition \eqref{eq 1 lem 3.5} provided that $M$ is large. Write $$G_j = \sum_{i = 1}^j g_i = \sum_{i = M}^{M+j}b_i,$$ Let $C$ be a sufficiently large absolute constant, and suppose that
	\begin{equation}\label{eq 3 lem 3.5}
		j - G_j \geq \max \big\{-1,(\min\{G_j,k-G_j\})^{1/7}-C\big\} \quad (1 \leq j \leq \Tilde{v}).
	\end{equation}
	For $b_j$'s chosen as described, we group the terms of our sum by the value of $m = G_i$. Since $G_i$ is non-decreasing and takes integer values between $0$ and $k$, we have
	\begin{align*}
		\sum_{j =M}^h 2^{-j +b_M+\cdots +b_j} &= 2^{1-M}\sum_{i=1}^{\Tilde{v}} 2^{-(i -G_{i})} \\
		&= 2^{1-M} \sum_{m=0}^k \sum_{\substack{1 \leq i \leq \Tilde{v} \\ G_i = m}} 2^{-(i-m)}.
	\end{align*}
	For each $m \in \{0, \dots, k\}$, let $I_m = \{i \in [1,\Tilde{v}] : G_i = m\}$. If $I_m$ is non-empty, let $i_m$ be its smallest index. As $i$ increases within $I_m$, the difference $i-m$ increases by $1$ at each step. Thus, the inner sum is bounded by a geometric series:
	$$ \sum_{i \in I_m} 2^{-(i-m)} \leq \sum_{\ell = 0}^\infty 2^{-(i_m - m + \ell)} = 2 \cdot 2^{-(i_m - m)}. $$
	By \eqref{eq 3 lem 3.5} evaluated at the starting index $i_m$, we know $i_m - m \geq (\min\{m,k-m\})^{1/7} - C$. Therefore,
	$$ 2 \cdot 2^{-(i_m - m)} \leq 2^{C+1} 2^{-(\min\{m,k-m\})^{1/7}}. $$
	Summing this bound over all possible values of $m$, we obtain
	\begin{align*}
		\sum_{j =M}^h 2^{-j +b_M+\cdots +b_j} &\leq 2^{2-M+C} \sum_{m=0}^k 2^{-(\min\{m,k-m\})^{1/7}} \\
		&\leq 2^{3-M+C} \sum_{m=0}^{\lfloor k/2 \rfloor} 2^{-m^{1/7}}.
	\end{align*}
	Finally, by comparing with an integral, we see that 
	$$\sum_{m=0}^{\lfloor k/2 \rfloor} 2^{-m^{1/7}} \leq \int_0^\infty 2^{-t^{1/7}}dt \leq 2^{20}.$$
	Hence, from \eqref{L to factorials} we conclude
	\begin{equation*}
		\sum_{\substack{a \in \mathscr{A}(\bm{b}) \\ \omega(a) = k}} \frac{L(a)}{a} \gg_{c,\kappa}  \frac{(2 \density \log 2)^k}{b_M! \cdots b_h!} \bigg[ 1 + \sum_{j =M}^h 2^{-j +b_M + \cdots +b_j} \bigg]^{-1} \ge \frac{1}{1+2^{23+C-M}} \cdot \frac{(2 \density \log 2)^k}{g_1! \cdots g_{\Tilde{v}}!}.
	\end{equation*}
	Since $M$ depends only on $c,\kappa$, we can rewrite this as
	\begin{equation}\label{L to factorials 2}
		\sum_{\substack{a \in \mathscr{A}(\bm{b}) \\ \omega(a) = k}} \frac{L(a)}{a} \gg_{c,\kappa} 2^{-C} \cdot \frac{(2 \density \log 2)^k}{g_1! \cdots g_{\Tilde{v}}!}.
	\end{equation}
	The quantity $(g_1! \cdots g_{\Tilde{v}}!)^{-1}$ is naturally interpreted as the area of a geometric region in $\R^k$. For $\bm{g} = (g_1, \dots, g_{\Tilde{v}})$. define
	$$\mathcal{R}(\bm{g}) = \{x \in \R^k : 0 \leq x_1 \leq \cdots \leq x_k <v , x_{G_{i-1}+1},\dots,x_{G_{i}} \in [i-1,i)\;1 \leq i \leq \Tilde{v}\},$$
	with the convention that $G_0 = 0$. We then have
	\begin{equation} \label{eq 3.5 lem 3.5}
		\frac{1}{g_1! \cdots g_{\Tilde{v}}!} = \Vol(\mathcal{R}(\bm{g})).
	\end{equation}
	Let $$\mathscr{G} = \Big\{\bm{g}=(g_1, \dots, g_{\Tilde{v}}): \sum_{i = 1}^{\Tilde{v}} g_i =k, \text{ and } \eqref{eq 2 lem 3.5}, \eqref{eq 3 lem 3.5}\text{ hold} \Big\}.$$
	We will show that
	\begin{equation}\label{eq 4 lem 3.5}
		\Tilde{v} \mathcal{Y}_k (\Tilde{v},C) \subseteq \bigcup_{\bm{g} \in \mathscr{G}} \mathcal{R}(\bm{g}).
	\end{equation}
	Given $\xi_i \in \mathcal{Y}_k(\Tilde{v},C)$, we set $x_i = \Tilde{v}\xi_i$. We claim that $x_i \in \mathcal{R}(\bm{g}),$ where $g_i$ is defined to be the number of $x_j \in [i-1,i)$. Let us verify that \eqref{eq 3 lem 3.5} holds for such an $x_i$. By the definition of $g_i$, we have $j > x_{G_j}$ for each $j$. Combining this with property (iii) of $\mathcal{Y}_k(\Tilde{v},C)$, we see that
	$$j > x_{G_j} \geq \max\{G_j-1, G_j + (\min\{G_j,k-G_j\})^{1/7} -C\}.$$
	Hence, $j - G_j \geq \max\{-1,(\min\{G_j,k-G_j\})^{1/7}-C\},$
	and so \eqref{eq 3 lem 3.5} is satisfied. Similarly, by condition (ii) in the definition of $\mathcal{Y}_k(v',C)$, we have $x_{M+i^2} > i,$ so $g_i \leq M+i^2$ for $(1 \leq i \leq \sqrt{k-M})$. One can similarly show by using the condition $x_{k-(M+i^2)+1} < \Tilde{v}-i$ that $g_{\Tilde{v}-i+1} \leq M+i^2.$ Hence, \eqref{eq 2 lem 3.5} is satisfied. 
	
	 This shows that $\bm{x} = (x_1,\dots, x_k) \in \mathcal{R}(\bm{g})$ for a certain $\bm{g} \in \mathscr{G}$, and so \eqref{eq 4 lem 3.5} holds. As such, we have
	\begin{equation}\label{eq 5 lem 3.5}
		\Vol\Big(\bigcup_{\bm{g} \in \mathscr{G}} \mathcal{R}(\bm{g})\Big) \ge \Tilde{v}^k \Vol\Big( \mathcal{Y}_k(\Tilde{v},C)\Big).
	\end{equation}
	Combining \eqref{L to factorials 2}, \eqref{eq 3.5 lem 3.5} and \eqref{eq 5 lem 3.5}, we have proven that
	$$\sum_{\substack{a \in \mathscr{A}(\bm{b}) \\ \omega(a) = k}} \frac{L(a)}{a} \gg_{c,\kappa} (2 \Tilde{v} \density \log 2)^k\frac{\Vol\big( \mathcal{Y}_k(\Tilde{v},C)\big)}{2^C}.$$
	In order to complete the proof, we must replace $\Tilde{v}$ with $v$ in our final estimate. To do so, simply note that 
	$$\Tilde{v}^k = v^k (1-(2M)/v)^k \geq e^{-2M} \gg_M v^k.$$ 
	Since $M$ depends entirely on $c,\kappa,$ we can replace $\gg_M$ with $\gg_{c,\kappa}$. This completes the proof.
\end{proof}

\end{document}
